\section{Assembly for exponents greater than two}
\label{ha:section}\label{sec:high-assembly}

We now choose the data in the carrier, quantizer, wall, and mesh
constructions simultaneously.  Throughout this section, put
\[
 \Lambda=\Omega',\qquad E=\int_\Omega |Df|^p,\qquad p>2.
\]
Matrix norms without a subscript are Frobenius norms.  All local
finiteness statements refer to the open source or target, rather than
to their closures.

Our objective is to prove Theorem~\ref{main:theorem} by constructing
an approximant with any prescribed positive error in
Equation~\eqref{main:convergence}.  The construction first gives a
locally bi-Lipschitz homeomorphism.  The last use of
Theorem~\ref{sm:smoothing} gives the smooth map.
We keep separate the constants allowed to multiply an
initial energy tail and those allowed only in later absolute-error
prescriptions.  This distinction is essential: the constants of an
individual topological normalization or an individual chart need not
be bounded in terms of $p$.

Two kinds of estimates enter the proof.  Ordinary meshes give the
\emph{weak estimate}: an energy bound in terms of the carrier energy.
Around the retained rank-one and rank-two data, calibrated meshes give
the \emph{sharp estimate}: an energy bound close to that of the affine
comparison.  Label pinning controls the mean gradient there, so a
quantitative uniform-convexity argument yields strong gradient accuracy;
this is the classical energy-to-strong-convergence principle associated
with Clarkson~\cite{Clarkson1936}.  The precise pointwise estimate used
here is proved in Lemma~\ref{ha:saturation}.  The complete order of choices
is collected in Remark~\ref{ha:scale-order}.

The estimates have four spatial roles.  The sharp estimate recovers
the derivative on small source patches containing almost all the
retained rank-one and rank-two energy.  Reserved full-rank cubes
will instead receive their affine comparison maps at the end.
When $2<p<3$, marked interiors are also set aside, for a later
source compression.  On the remaining region, and on the shells
and collars needed to make these replacements, the weak estimate
must give small energy.  These buffered regions can overlap;
their separate bounds will be imposed simultaneously.  We first
reserve the full-rank cubes and the carrier budgets, then choose
the sharp patches and realize the wall parameters that recover
their gradients.

\subsection{Initial parameters and carrier budgets}
\label{ha:initial}

Choose a small aspect $\gamma>0$ for the protected rank-one rectangles;
use squares for rank two.  Choose a small slot-loss parameter $\eta>0$,
and then $\zeta,c_*>0$ sufficiently small compared with $\eta$.
Here $\zeta$ controls occupied root length and sampling density, and
$c_*$ is the ratio of a true target interval length to its root weight.
The ordinary weak constant, denoted $C_{\mathrm w}$, may depend on
\[
                       p,\gamma,\eta,\zeta,c_* ,
\]
but it is independent of all later jet truncations, chart bounds,
mesh resolutions, and guard spacings.  The physical bulk meshes and
clustered templates of Proposition~\ref{ms:prestacks} supply precisely
this kind of constant.  Protected rectangle aspects are included in
$C_{\mathrm w}$.  We will prove a leading sharp error bound
\begin{equation}\label{ha:leading-error}
 C_p(\eta+\gamma)(1+E)+\text{freely small errors},
\end{equation}
where $C_p$ is independent of $\gamma,\eta$ and of the finite jet
bounds.  Thus $\gamma$ and $\eta$ can be fixed at the outset to make
the first term as small as required.

In what follows, a freely small error is prescribed only after every
factor that will multiply it has been fixed.  On noncompact parts we
use a locally finite compact cover and positive prescriptions with
sum as small as required.  A finite number of additional local weights
can be included in these prescriptions by dividing the allowed error
by their maximum.  Equivalently, for a countable collection of local
costs, prescribe the $j$th weighted cost to be less than
$\delta 2^{-j}$, after its weights are known.  Positive tolerance minorants and simultaneous locally summable
choices are supplied by Lemma~\ref{pre:local-tolerances}.  This
convention never replaces the initial choice of tails against
$C_{\mathrm w}$.

By Lemma~\ref{hp:regularity}, we may choose compact sets of regular
rank-one and rank-two points covering those ranks except for an
arbitrarily small tail in the measure
\[
                         d\mu=(1+|Df|^p)\,dx.
\]
First make this tail small compared with the already fixed weak
constant.  On the selected sets the positive singular values can be
bounded above and bounded away from zero.  Separately choose a compact
set of full-rank points with finite jet and inverse-jet bounds,
leaving an equally affordable tail in $\int |Df|^p$ on that rank.
The full-rank points and the deficient compact data are separated.

Use Proposition~\ref{hp:full-rank-correction} to reserve finitely many
full-rank patches.  Write $B_i$ for their inner cubes and $b_i$ for
the positive affine comparison jets.  The enlarged support and
comparison cubes are disjoint across distinct patches.
In normalized cube coordinates,
let $t_i=0$ on $\partial B_i$ and let $s>0$ be the cut-band parameter.
The finite jet factors are fixed before $s$ is made small.  For any
fixed large $C'$, all bands
\begin{equation}\label{ha:full-bands}
                         |t_i|\le C's
\end{equation}
can consequently have arbitrarily small total $|Df|^p$ integral,
even with the jet factors and $C_{\mathrm w}$ included.  The constant
$C'$ includes the finitely many buffer enlargements used below.
The power-mean and normalized value tests at the full-rank points
are then imposed on sufficiently small outer cubes.  Vitali selection
and finite truncation give coverage by the inner cubes, and even by
$\{t_i<-3s\}$, up to the required full-rank tail.  We retain
\begin{equation}\label{ha:affine-full-tests}
 \sum_i\int_{B_i}|Df-Db_i|^p\ \text{as small as required},
 \qquad
 |b_i^{-1}f-\mathrm{Id}|\ll cs
 \quad\text{in normalized cube units}.
\end{equation}
The cube sizes also make the eventual value cost of these corrections
arbitrarily small.  Subsequent value approximations will preserve
the second test with room.

Before fixing the flattening and endpoint data, choose a positive
continuous source tolerance $a$ small enough for the desired value
error and the finite normalized full-rank tests, reserving room for
all later changes.  In Proposition~\ref{hp:flattening} use a target
tolerance $a_\Lambda(y)\le a(f^{-1}(y))/4$.  Then
$|F_\kappa-f|<a/4$ for every $0\le\kappa\le1$.  After this share
has been reserved, the two carrier rounds use the remaining
allowance in $a$, as in Proposition~\ref{hp:carrier}.
When $2<p<3$, fix the marked-collar enlargement factor large enough
for the fixed two-sided buffers used below, before assigning the
carrier's facet budgets; only the collar widths are chosen after
both rounds.

Apply the flattening and two-round carrier construction, using
Proposition~\ref{hp:flattening} and Proposition~\ref{hp:carrier}, and then
the quantizer calibration of Lemma~\ref{qt:calibration}.
Choose the line directions and the flattening tolerance $\lambda$
so accurately that, on the ultimately retained rank-one data, the
range direction of $Dq_0$ differs from the long rectangle axis by
$o(\gamma^2)$.  The derivative of the flattening has relative error
$C\lambda$ there.  The central fraction of the quantizer uses a common
scalar times orthogonal projection onto the protected plane; its
scalar $1/a_*$ can be taken arbitrarily close to one.  Trimming to
these fractions loses arbitrarily little $\mu$-mass.  Kernel inclusion
follows from Lipschitz composition on the retained data, and the
positive singular-value bounds preserve the rank when the errors are
small enough.

Denote by $K_c$ the common compact data surviving the carrier rounds,
before the later phase losses.  In both rounds the implant supports
\emph{and their regular comparison boxes} avoid $K_c$ with room.
Indeed the regular assignments are first restricted to compact sets
separated from the central data; only finitely many operations meet
a neighborhood of $K_c$.  The two rounds can therefore use common
positive separations.  In particular,
\begin{equation}\label{ha:q0-weak}
 F_b=F_\kappa\quad\hbox{near }K_c,
 \qquad q_0=TJF_b,
 \qquad |Dq_0|\le C|DF_b|\quad\hbox{a.e.}
\end{equation}
All supports and comparisons are chosen sufficiently fine relative
to the reserved full-rank patches and their bands.  The regional
charging conclusion of Proposition~\ref{hp:carrier} has only the two
preassigned buffer enlargements, so this is a requirement on those
comparisons before they are fixed.

Here are the resulting regional budgets, stated explicitly for their
later uses.  Put
\[
                     Z=\Omega\setminus\bigcup_i\{t_i\le-s\}.
\]
Outside $K_c$ and the marked cubes, the $|DF_b|^p$ integral on $Z$ is
small against $C_{\mathrm w}$.  Each regular implant contributing
there charges a disjoint comparison away from $K_c$.  Exterior
comparisons in the second round also avoid the first-round marked
cubes.  Both charging enlargements stay in $\{t_i\ge-3s\}$ for every
$i$.  The unchanged derivative is bounded by $C|Df|$, while exceptional
implant costs are confined to marked interiors.  The rank-zero set
contributes no $|Df|^p$ cost.  The same assertion holds on the enlarged
full-rank shells outside marked interiors, with the finite correction
jet factors included, because their charging enlargements lie in the
larger bands \eqref{ha:full-bands}.

If $2<p<3$, use the marked cubes $U$ of
Proposition~\ref{hp:carrier}.  Their enlarged union has arbitrarily
small $\int(1+|Df|^p)$.  After both carrier rounds choose collar widths
$d_U>0$, with fixed-factor two-sided buffers and mutual separation,
so that
\begin{equation}\label{ha:collar-budget}
 \sum_U\frac{\ell_U}{d_U}
             \int_{\mathcal C_U}|DF_b|^p
                  \quad\hbox{is as small as required}.
\end{equation}
Here $\ell_U$ is the side scale, and $\mathcal C_U$ is an enlarged
collar of width comparable to $d_U$ about $\partial U$.
The prescribed smallness includes $C_{\mathrm w}$ and the full-rank
jet factors wherever those tests overlap.  The unchanged-base slice
estimate in Proposition~\ref{hp:carrier} allows the widths to be
chosen after the rounds, however small the unchanged neighborhoods
have become.  The marked cubes, their enlargements and collars avoid
the deficient protected data.  Their sizes are also fine enough to
make the oscillation of $f$ on each enlarged cube meet all required
value tests.  There are no marked cubes for $p\ge3$.  The weights
$\ell_U/d_U$ are now fixed for all later absolute-error prescriptions.

\subsection{Polar localization and the active derivative estimate}
\label{ha:active-localization}

The central-tube modification will change $q_0$ to a map $q_*$
with only finite local Sobolev bounds in its low-radius regions.
We arrange that the final nonstrict map is constant there.
Thus the weak derivative estimate will use $q_*=q_0$ wherever
an output parameter is active; the larger low-region bounds will
enter only the crude capacities needed to pin the labels.

We next choose the protected arrays, phases and subcells from
Lemma~\ref{qt:calibration} and the polar setup of
Lemma~\ref{qt:polar}.  Arrays can be arbitrarily fine for the value
accuracy and for the quantizer's maximum-activity and clearance tests.
Phase selection retains the protected $\mu$-weight in rectangular
interiors, up to arbitrarily small loss, and avoids atoms at the
centers and the rays to corners.  Additional subedge divisions along
straight protected sides can also avoid ray levels of positive
$\mu$-mass: before overlaying the polygon meshes, take strict interior
cuts in the individual clipped diagrams and slide these finitely
localized cuts generically in their available intervals.

All full-cell and lower-face geometries, and their compact-local
aspect bounds, are now fixed.  For unprotected polygons, use the
original-vertex localization in Lemma~\ref{qt:subdivision} before
fixing the fine subcells.  In detail, the worst relevant aspect factor
on each true polygon is finite at this stage.  Moreover $Dq_0=0$ a.e.
on each vertex level set.  Thus the integral of $|Dq_0|^p$, multiplied
by that fixed aspect factor and any already assigned local weights,
tends to zero on the preimages of shrinking vertex neighborhoods.
Properness makes these tests locally finite.  Choose the neighborhoods
with summable costs, and then fit the fine subcells.  Everywhere else
the aspect factors in the ordinary weak estimate are absolute or
are the already fixed protected rectangle factor.

Choose first the conical boundary-layer widths in the full output
cells, as in \eqref{wl:conical-extension}.  Their weighted excess
costs can be made summably small before choosing the central activation
radii.  To see the order, in one boundary polygon the two-factor
interpolation has a Lipschitz bound depending on its fixed aspect and
$c_*$: its radial factor satisfies $\rho/r\le C$ and
$|\rho'|\le C/c_*$, while the inverse angular derivative contributes
an aspect factor times $1/r$.  The $r$ cancels.  Once the layer
thickness is fixed, its normal derivative is controlled by requiring
the factor motions to be sufficiently small.  The parametrizations
of the full input blocks are already fixed exact locally bi-Lipschitz
data, so these are ordinary absolute-continuity choices on known
maps.

Now take the activation radii $a_D$ sufficiently small for those
motion tests, and retain compact data of both positive deficient
ranks where $r>5a_D$, strictly inside smooth sectors.  Center and ray
avoidance makes the additional loss arbitrarily small; no absolute
continuity of the projected protected measure is being assumed.
Choose $r_{u,D}\ll a_D$ and use
Lemma~\ref{qt:central-carrier} to obtain
$M_y,h_*,M_x,F_*,q_*$.  The exterior changes from $q_0$ are confined
to deep low regions within individual subcells, which we may require
to satisfy
\begin{equation}\label{ha:deep-changes}
                         r(q_*)<a_D/100.
\end{equation}
The relative radius is set equal to one on subedges, giving a
continuous function on the two-complex.  All openings, central changes
and labels retain the preliminary fine value tests.

Write $K^\sharp\subset K_c$ for the final protected compact set.
It lies in the interior of $M_x$; the slab heights, tube clearances
and strict sector insets give room to require $q_*=q_0$ on its
neighborhood.  Lemma~\ref{qt:calibration} gives, a.e. there,
\begin{equation}\label{ha:relative-calibration}
                   |Dq_0-Df|\le C(\lambda+1-a_*)|Df|.
\end{equation}
All losses from $K_c$ made so far are small against $C_{\mathrm w}$.

The wall barriers will require arbitrarily fine local closeness of
$q_h=Th$ to $q_*$.  Prescribe that closeness so that the relevant radii
are comparable, including on all mesh-neighborhood enlargements.
The radii are positive and locally bounded away from zero on $M_x$.
Whenever $q_h$ lies on an edge, or in the possible active face range
$r(q_h)\ge a_D/2$, the comparison neighborhood is contained in an
open high-budget region where $q_*=q_0$.  Such neighborhoods exist by
\eqref{ha:deep-changes}.  Prescribe sufficiently small stair shifts,
also relative to $L_D$ in the logarithmic coordinate, that the
nonstrict map is constant below this active range.

Let $G$ be the upper physical pre-parameter slope supplied by the
mesh and wall constructions, allowing a fixed factor for sums of
two slopes.  On the active region, Lemma~\ref{wl:stairs},
\eqref{wl:target-speeds}, and Proposition~\ref{wl:parameter-cost} give
\begin{equation}\label{ha:active-bound}
                      |D(P_0q_h)|\le Cr(q_h)G
                      \quad\hbox{a.e.}
\end{equation}
Indeed a nonzero stair derivative occurs in an occupied stack.  Its
root parameter has one pre-parameter derivative or a sum/difference
of two derivatives on a switched diagonal.  Every intervening
OUT-to-IN identification has absolute slope one, so there is no
factor depending on itinerary length.  The occupied stair has root
speed at most $1+O(\zeta)$, and the forward polar map costs $O(r)$,
including the activation ramp.  On edge interiors this is the
perimeter estimate with $r=1$; on vertex levels the derivative
vanishes.  The formulas suffice a.e. at knots and piece interfaces:
$h$ is locally bi-Lipschitz, the face and edge product projections
are absolutely continuous for their respective dimensional measures,
and derivatives vanish on constant level sets.  In full blocks,
away from the already paid conical layers, $P_0q_h=q_0$.  On the
remaining central tube interiors outside $M_x$ the nonstrict map is
constant by the exact correspondence.  Thus only the prescribed
high-budget regions require the weak estimate
\eqref{ha:active-bound}.

\subsection{Calibrated source patches}
\label{ha:patches}

Use the fixed cuboidal and coarse chart structures on $M_x$.
Their walls and nonsmooth piece loci are locally finite ambient-null
sets.  Trim $K^\sharp$ away from them with room, losing arbitrarily
little $\mu$-mass.  At a retained Lebesgue gradient point $z$, within
a single smooth $(D,e)$ sector, put
\begin{equation}\label{ha:Y0}
              Y_0=(t(q_0),s(q_0)),\qquad B=D_xY_0(z).
\end{equation}
We choose a nonsingular spatial basis $V$ as follows.  If $B$ has
rank two, let $k$ be a unit vector in its kernel and take
\[
                    V=[B^\dagger,k/\|B\|_{\mathrm{op}}].
\]
The columns of $B^\dagger$ are orthogonal to $k$, and the singular
values of $V$ show that
\begin{equation}\label{ha:rank-two-basis}
             BV=[I_2,0],\qquad
             \|V^{-1}\|_{\mathrm{op}}=\|B\|_{\mathrm{op}}.
\end{equation}
If $B$ has rank one, write $B=a v^T$, where $|v|=1$, put
$\alpha=|a_1|+|a_2|$, and choose an orthogonal $Q$ with $Qe_1=v$.
Then $V=Q/\alpha$ gives
\begin{equation}\label{ha:rank-one-basis}
 BV=\begin{pmatrix}b_1&0&0\\b_2&0&0\end{pmatrix},
 \qquad |b_1|+|b_2|=1.
\end{equation}
For the rank-one rectangle,
\begin{equation}\label{ha:b-small}
                        \min(|b_1|,|b_2|)\le C\gamma.
\end{equation}
On a long-side sector the long-axis direction has only perimeter
speed, of order $1/r$.  On an end sector it has logarithmic speed
of order $1/r$ and perimeter speed at most $C\gamma/r$.  The
$o(\gamma^2)$ direction error, combined with inverse sector cost
$O(1/(\gamma r))$, preserves these estimates.  For rank two the
rectangles are squares.  In both cases Lemma~\ref{qt:polar} and
\eqref{ha:relative-calibration} give the absolute bound
\begin{equation}\label{ha:V-cancellation}
                r(z)\|V^{-1}\|_{\mathrm{op}}\le C|Df(z)|,
                \qquad r(z)=r(q_0(z)).
\end{equation}
This is the cancellation needed when the normalized sharp estimate
is returned to physical coordinates: the target reconstruction
costs $O(r)$, while the source change of coordinates costs
$\|V^{-1}\|_{\mathrm{op}}$.  Their product is controlled by the
original gradient, without a factor from the finite jet bounds.
On the retained compact set, $V$ and $V^{-1}$ have finite bounds:
use the positive rank, the singular-value truncations, the fixed
sector insets, and the finite chart bounds.

In a smooth coarse chart $x=g(\xi)$ use aligned coordinates
\[
 \xi=\xi_z+A y,\qquad A=Dg(\xi_z)^{-1}V.
\]
Choose small outer $y$-cubes centered at $y=0$, corresponding to $z$,
with nested inner cubes.
On their enlargements require: a single smooth sector with small
label oscillation; chart derivatives close to their center values;
and arbitrarily small power means of
\[
                         DY_0-B,\qquad Df-Df(z).
\]
All required eccentricities are bounded on the retained compact, so
these are a Vitali differentiation family.  Select disjoint outer
cubes and then a finite subfamily.  Choose the nested ratios close
enough to one that the inner cubes still capture the required weight,
including a further inset for residual tests.  We call them the
\emph{saturation cubes}.  Their physical images are separated from the
boundary of $M_x$, the marked cubes and the full-rank correction
supports.

There is no circular dependence between their small tests and their
geometric constants.  First fix the finite chart, sector, basis,
slot-shape and template bounds on the prospective compact data.
The finite triangulation lists in Lemma~\ref{ms:snapping}
depend on bounded cardinalities and fixed slot margins, not on
microscopic noncluster angles or master widths.  The transition mesh
bounds in Lemma~\ref{ms:meshes} depend on the fixed bases, not on the
relative thickness of patch buffers.  Next fix the level-edge
count tolerance, then smaller chart-error and jitter fractions,
and then still smaller power-mean data errors, allowing the inverse
jitter-density factors.  Finally choose patch radii realizing these
tests.  The mesh step $H$ and then $W/H$ can be arbitrarily smaller.

Thin patch margins, including equal-weight, clustering and Delaunay
exit belts, are chosen after the fixed local derivative and density
factors, so their volume and needed Sobolev integrals are small.
They can be resolved without deterioration of the shape factors.
Keep every equal-weight convolution cost inside enlarged pure lattice
regions where the affine tests remain valid.  Microscopic master
widths occur only inside the snapping construction; leave several
layers of widths comparable to $s_\sigma$ in the clustered
equal-weight belt before returning to ordinary widths.  Sharp and
mixed-position tetrahedra use the noncluster generic convention;
there is no robust-template motion of an opponent system in the
sharp bulk.  The wholly clustered convention resumes past the
buffer.  Outside-chart and unequal-orientation transitions use the
fixed coarse-shape weak bound on their paid regions.

\subsection{Ambient edge tests and weak costs}
\label{ha:weak-costs}

Outside saturation interiors use the physical-mode overrides of
Lemma~\ref{ms:meshes} on smooth coarse chart interiors.  Reserve small
neighborhoods of coarse facets, chart-piece seams, patch and Delaunay
transitions, and nonphysical mesh regions.  Their integrals of
$1+|Dq_*|^p$, with every assigned compact-local factor and residual
weight, can be made summably small.  Indeed these factors are fixed
before the neighborhood widths, and finitely many inset physical
patches cover a coarse compact except for arbitrarily small weighted
error.  Enlarge the error sets slightly to include neighboring mesh
stars.  The actual cut-boundary adjustment and exact-region jitter
taper can be made equally thin, later if necessary.

The last assertion uses the width-independent bound in
Lemma~\ref{ms:boundary-collar}.  The horizontal warp and chord inverse
response have fixed bounds, and motions depending on $d/d_0$ have
size comparable to the boundary mesh scale, which is much smaller
than $d_0$.  The slopes of $d_0$ are controlled.  Thus the label
slope, including its conversion to a root parameter, has no inverse
power of $d_0$; a fixed factor $1/c_*$ is allowed.  Equal subdivision
of the boundary chords changes only displacement differences in
tangential derivatives and adds no subdivision factor.

The exterior residual test lies in $Z$, outside slightly shrunken
saturation interiors, and enters a marked interior only through
its enlarged collar.  The other two tests are the buffered
full-rank correction shells and the marked collars with weight
$\ell_U/d_U$.  Mesh stars and jitter displacements will fit inside
the enlargements already reserved for these tests.

\begin{lemma}[Ambient edge tests and weak costs]\label{ha:ambient-tests}
With the preceding geometry and local weights fixed, the ordinary
weak costs are controlled by the carrier budgets with constant
$C_{\mathrm w}$, plus freely small absolute errors.  The control is
uniform over meshes sufficiently fine for prescribed buffer and
guard requirements.  It applies simultaneously to the exterior
residual, the full-rank correction shells, and the marked collars
weighted by $\ell_U/d_U$.
\end{lemma}

\begin{proof}
For a mesh edge or edge portion $l$, let $W_l$ be the sum of the
weights of its initial central hits.  On whole edges these sums
are upper bounds for the normalized systems: cleaning, opening and
normalization do not increase any individual wall--edge count.
Edge portions are used only to estimate the initial counts by
coarea; their contributions are added before normalization.
On an edge portion outside the modified
boundary collar, conditional on a suitable ACL mesh,
Lemma~\ref{wl:stairs} gives
\begin{equation}\label{ha:weighted-coarea}
 W_l\le(1+C(\zeta+c_*))\operatorname{Var}_l(q_*)+\varepsilon_l.
\end{equation}
Here the variation is the sum of star variation for the radial family
and graph arclength variation for the perimeter family; the two may
also be tested separately.  The additive errors $\varepsilon_l$ may
be as small as required relative to the actual edge and its weights.
Each feature has only finitely many relevant edge tests on a compact
localization, by properness.

The coefficient close to one is the weighted sampling density of
Lemma~\ref{wl:stairs}, not the width of $I_j$.  Its stratified
sampling argument, applied to the integrable crossing-count tests
in Lemma~\ref{ms:edge-tests}, gives the inequality with arbitrarily
high probability after the edges have been chosen.  It also permits
a common weight on finitely many features and all almost-sure null
avoidances.

In a high-budget region, the two graph speeds satisfy
\begin{equation}\label{ha:graph-speed}
                 |D(\hbox{graph labels})|
                    \le \frac{C_{\mathrm{asp}}}{r}|Dq_0|
\end{equation}
along the relevant ACL edge paths.  The bound is a stratumwise length
bound; an inverse ambient dihedral angle is unnecessary.  On a radial
spoke the perimeter-vertex coordinate is constant, and
$q=d_D+r(\xi-d_D)$ gives
$|dt|\le (L_D/|\xi-d_D|)|dq|/r$.  On a subedge the perimeter
projection is the identity and the star coordinate is constant.
At vertices or graph nodes the corresponding coordinate speed
vanishes a.e. on its level set.  Restrictions at density and
differentiability times have the ambient ACL derivatives, so these
bounds combine along the path.  Local Lipschitz graph coordinates,
possibly with larger fixed seam constants, ensure the required
absolute continuity.  The estimate itself uses just the sector
aspect factors.

On modified cut-collar portions replace \eqref{ha:graph-speed} by
the fixed local label-slope bound of
Lemma~\ref{ms:boundary-collar}.  The moving graphs and chords have
inverse-response bounds on the allowed bins, also across the
triangle pieces.  Restricted one-dimensional coarea on the disjoint
allowed parameter strips gives the corresponding weighted-count
bound in terms of these prepared labels, with the stated fixed
factors.  The central prepared trace family is fixed before sampling;
its later interval-ribbon adaptation keeps the sampled central
chords unchanged.  We split an edge into its unchanged and modified
collar portions when applying these bounds.  The depth-stretched
portions use the adjoining exact-region bounds.

For an ordinary all-clustered tetrahedron $\sigma$ of physical scale
$h_\sigma$, Proposition~\ref{ms:prestacks} and
\eqref{ha:active-bound} give the active slope majorant
\begin{equation}\label{ha:weak-tet-slope}
                 C r_\sigma h_\sigma^{-1}
                     \max_{l\subset\sigma}W_l.
\end{equation}
The radii on the required stars are comparable to $r_\sigma$.
There is no fixed baseline in \eqref{ha:weak-tet-slope}: a
contributing pre-parameter has width proportional to its root
weight, and zero crossings require no disk.  Apply
\eqref{ha:weighted-coarea}, \eqref{ha:graph-speed} and H\"older's
inequality on an edge $x_l(u)$, $0\le u\le1$, parametrized at speed
$O(h_\sigma)$.  The $p$th power of
\eqref{ha:weak-tet-slope}, times $|\sigma|$, is at most an arbitrarily
small additive error plus
\begin{equation}\label{ha:edge-power}
 C h_\sigma^3\sum_{l\subset\sigma}
       \int_0^1 C_{\mathrm{asp}}(x_l(u))^p
                    |Dq_0(x_l(u))|^p\,du.
\end{equation}
Only tetrahedra in deterministic high-budget buffers require this
estimate.  Costs which become inactive after pinning the labels are
not included in the weak objective.

Lemma~\ref{ms:edge-tests} supplies the ambient expectation estimate:
on an unconstrained jittered edge, the point at fixed normalized
edge time has density at most $C h_\sigma^{-3}$ in a neighboring
star, and the normalized edge speed is at most $C h_\sigma$.
The lemma also supplies ACL and the chain rule.  Thus for any
nonnegative ambient integrable weight $F$,
\begin{equation}\label{ha:jitter-expectation}
 \mathbb E\left[h_\sigma^3\int_0^1F(x_l(u))\,du\right]
       \le C\int_{\operatorname{star}^+(\sigma)}F(x)\,dx.
\end{equation}
Here the enlarged stars have bounded overlap.  Jitter radii are
chosen to stay within those stars.  The taper into nonrandom data is
wholly in exact regions, where fixed local Lipschitz bounds replace
the expectation estimate.  In the physical bulk, shapes, displacement
fractions and overlap are universal.  Keep any additional mesh layers
needed to reach that bulk inside the paid transition regions.
Summing \eqref{ha:jitter-expectation} therefore charges
\eqref{ha:edge-power} to the prescribed ambient budgets with only
$C_{\mathrm w}$ in the ordinary bulk.  In the exceptional zones it
charges $1+|Dq_*|^p$ with the already fixed local factors.  Near the
original polygon vertices it charges the affordable weighted
$|Dq_0|^p$ selected before subdivision.

We now charge the deterministic envelopes specified above.  The
part of the exterior-residual envelope in $K_c$ is small in
$|Df|^p$ by weighted capture and \eqref{ha:q0-weak}; on its complement
use the regional $DF_b$ budget and \eqref{ha:collar-budget}.  The
full-rank shells and marked-collar envelopes lie in their preassigned
enlargements and avoid the sharp construction.  Mesh stars and
jitter displacements are finer than all these buffers.  The sharp
transition and equal-weight belts use the margin and leakage estimate
proved below.  Add the already small full-cell conical-layer costs;
elsewhere in full blocks the map is $q_0$.  The objectives are a
bounded number of aggregate nonnegative costs, and their local
weighted summands can be combined into one objective.  Their
expectations have the asserted smallness.  Lemma~\ref{ha:capacities}
selects the mesh and sample data simultaneously with the crude guard
capacities; Proposition~\ref{ha:initial-walls} realizes the actual
wall systems with these same bounds.
\end{proof}

\subsection{Sharp counts, convolution and discarded regions}
\label{ha:sharp-errors}

On an enlarged calibrated patch let $H$ be the common aligned Kuhn
step.  The ambient volume jitter is included in its chart $g'$ from
$y$ to $x$.  The small first-chart errors and jitter fraction ensure
that $Dg'$ is as close to $V$ as prescribed, in relative matrix norm.
All relevant features use a common root weight $W$, with $W/H$
arbitrarily small.

\begin{lemma}[Sharp counts and discarded costs]\label{ha:sharp-counts}
The meshes and the subsequent stratified samples can be chosen so
that the sharp-template count conditions hold outside tetrahedra
of arbitrarily small aggregate pre-derivative $p$-cost.  More
precisely, for the two family counts on an edge $l$,
\begin{equation}\label{ha:sharp-edge}
 \frac{N_{l,i}W}{H}
 \le (1+C(\zeta+c_*))
          \left|(BV)_i\frac{\Delta y_l}{H}\right|+e_l,
 \qquad e_l\ge0,
\end{equation}
and the aggregate $\sum_lH^3e_l^p$ can be arbitrarily small after all
fixed template, transition and physical conversion factors.
The dirty portions left by Proposition~\ref{ms:sharp-contract} have
arbitrarily small aggregate evaluated pre-derivative $p$-cost in
aligned coordinates.
\end{lemma}

\begin{proof}
Use the separate coordinate variation tests in the single smooth
sector.  Along an edge, subtract the affine comparison $BVy$.
The resulting variation is bounded by the integral of
$|DY_0-B|$ times the bounded chart speed, plus the chart/jitter
derivative error times $\|B\|$.  Normalize by $H$ and use
H\"older's inequality.  The expectation estimate
\eqref{ha:jitter-expectation} bounds the sum of the resulting
$p$-errors by the power-mean data errors on the enlarged patch.
The inverse jitter-density factor may be large, but was fixed before
those tests.  The small chart/jitter terms are deterministic.
Conditional sampling gives the multiplicative factor in
\eqref{ha:sharp-edge} and arbitrarily small additive errors, which
are included in $e_l$.  There are only finitely many sharp patches,
and all matrix bounds are finite.  We may therefore impose an
aggregate error with any prescribed finite patchwise weights;
there is no requirement of simultaneous probabilistic control of a
relative mean on every individual small cube.

Fix first a count threshold much smaller than $\eta$, and also a
level-edge threshold $\delta'>0$, which may be arbitrarily small
after the $\eta$-dependent template constants.  Discard tetrahedra
where a relevant $e_l$ exceeds these thresholds.  On an inter-edge,
\eqref{ha:sharp-edge} bounds each independent count by
$(1+o(\eta))H/W$.  In parallel mode the sum of the two ideal counts
is $H/W$, by \eqref{ha:rank-one-basis}, so the combined excess is
also $o(\eta)H/W$.  On a level edge the ideal variation vanishes:
its count is at most $\delta'H/W$, with no multiplicative bias.
These distinct tolerances have distinct uses.  Inter-edge excess
controls fitting the lists into slots of spacing $(1-C\eta)W$;
level-edge errors control non-slabs and discrepancies in slab ranks.
Start the lists at common offsets of order $\eta H$ from the lower
axis end and perturb only within the unused slack.  Across a
regular tetrahedron the rank offsets of corresponding slabs differ
by at most $O(\delta'H/W)$, including the first band before the
second band in parallel mode.  These are the hypotheses of
Proposition~\ref{ms:sharp-contract}; no lower bound for the number
of surviving slabs is needed.

We give the cost calculation for discarded tetrahedra.  In the
finite enlarged equal-weight region, put
\[
 m_\sigma=1+\max_{l\subset\sigma}N_lW/H,
 \qquad
 M_\sigma=\sum_{\tau}\theta^{d(\sigma,\tau)}m_\tau,
\]
where $N_l=N_{l,1}+N_{l,2}$ is the total two-colour upper count,
the star graph has bounded degree, and $\theta$ is smaller
than a fixed reciprocal of that degree.  The initial weighted
counts remain upper bounds for all normalized counts, so they can
be used in these majorants.  Equation~\eqref{ha:sharp-edge} implies
$m_\sigma\le m_0+C e_\sigma$, where $m_0$ is an absolute baseline
and $e_\sigma$ is the largest of its finitely many edge errors.
The kernel has uniformly bounded row and column sums.  Thus its
convolution is bounded on every $\ell^p$, and
\begin{equation}\label{ha:convolution-error}
 M_\sigma\le M_0+L_\sigma,
 \qquad
 \sum_\sigma H^3L_\sigma^p
       \le C\sum_\sigma H^3e_\sigma^p.
\end{equation}
Adjacent $M$ values also have bounded ratios.

For a fixed threshold $\tau>0$, the number-weighted volume of
$\{e_\sigma>\tau\}$ is at most
$\tau^{-p}\sum H^3e_\sigma^p$.  Its $M^p$ cost is small as well:
pay the baseline $M_0^p$ by that volume and the remainder by
\eqref{ha:convolution-error}.  On $\{M_\sigma>2M_0\}$, the full
$M^p$ cost is bounded by a constant times the $L^p$ cost.  Fixed
neighbor layers of these sets have the same conclusions, using
bounded degree and adjacent ratios.  In thin patch margins within
the convolution region, pay the baseline by volume and the leakage
by \eqref{ha:convolution-error}.  Mixed-width tetrahedra in the last
quantitative belt can use both this majorant and the ordinary edge
bound; they still lie in the paid margin.

On low-leakage regular tetrahedra, fix the buffer multiple $N$
sufficiently large after the slot margins.  It accommodates the
flat translations, the auxiliary-cell trimming for snapping, and
the slow broadening to strip half-width
$a_1=(1-O(\eta))W/2$.  Proposition~\ref{ms:sharp-contract} gives
relative dirty volume at most
\begin{equation}\label{ha:dirty-fraction}
                           C_N(W/H+\delta').
\end{equation}
The pre-slope there is at most $CM_\sigma$, and $M_\sigma$ is
bounded on these tetrahedra.  First make $\delta'$ sufficiently
small, then $W/H$ sufficiently small.  The remaining bad-tetrahedron
and high-leakage costs are small by
\eqref{ha:convolution-error}, and the patch-boundary buffers are
paid margins.  Altogether the integral of the pre-slope to the
$p$th power on the dirty portions is arbitrarily small.  The
opposite-band gap and slope errors away from these portions are
$o(\eta)$ by the same level tolerance, negligible generic motions,
and $W/H$.  All constants used here were fixed before the final
aggregate error was prescribed.
\end{proof}

\subsection{Capacities before guards and simultaneous choices}
\label{ha:guard-schedule}

Here a capacity means a finite upper bound for a local integral
$\int G^p$, not a Sobolev capacity of a set.  The preceding estimates
are uniform over sufficiently fine mesh
resolutions.  They therefore allow the crude $G$ capacities to be
fixed before the guard nets, even though the final meshes must be
finer than the guard buffers.

There are two tasks.  We first select edges and samples for which
the count-based slope bounds fit the chosen capacities and energy
budgets.  We then construct the actual initial walls and pre-stacks,
retaining those very counts.  The second task is carried out in
Proposition~\ref{ha:initial-walls}; until then, the selected bounds
are upper budgets for the eventual slope density $G$.

\begin{lemma}[Prior capacities and simultaneous choice]\label{ha:capacities}
For each member of a locally finite family of buffered compact
barrier charts one can fix a finite capacity $A_j$, independently of
guard spacing and final mesh resolution.  After fixing these
capacities, one can choose the guards, mesh and samples so that
substitution of the selected edge counts into the slope bounds of
Proposition~\ref{ms:prestacks} gives local integral upper budgets
at most $A_j$.  The same data meet the weak aggregate bounds of
Lemma~\ref{ha:ambient-tests} and the sharp count and discarded-cost
bounds of Lemma~\ref{ha:sharp-counts}.
\end{lemma}

\begin{proof}
\emph{Uniform preliminary budgets.}
For the crude bound on a compact chart, include the upper bound for
every pre-slope, including ones that will later be inactive.
In nonsharp tetrahedra,
Proposition~\ref{ms:prestacks} bounds the slope by the weighted edge
variations with the fixed local chart, aspect and inverse-radius
factors, together with the allowed fixed width terms and the
cut-collar bounds.  Equations~\eqref{ha:weighted-coarea} and
\eqref{ha:jitter-expectation} give a finite expectation majorant
using $1+|Dq_*|^p$ on an enlarged compact.  The local factors are
bounded there.  In equal-weight regions the aligned pre-slope is
at most $CM_\sigma$, and the bounded convolution estimate charges
the whole enlarged patch to its $m$ budget.  Physical conversion
has only fixed compact factors.  Sharp overrides occur in a finite
interior compact.  Thus, for the $j$th barrier chart, there is a
finite number $B_j$ bounding the expected crude upper budget, uniformly
over all sufficiently fine resolutions and all permitted guard
gaps.

This expected cost is an upper-budget construction.  First require
the weighted sample counts to be bounded by their variations plus
the prescribed errors, conditional on the edges.  Take the $p$th
power of those variation upper bounds and then the ambient jitter
expectation.  No high moment of the unweighted crossing count is
being assumed.  The additive count errors can be chosen to obey
both the small aggregate objectives and the crude capacities.
The boundary-collar constants are independent of the excluded belt
widths and of $d_0$, and the mesh constants are independent of guard
spacing.  A guard-induced allowed component may be very short;
only the eventual sample weight, not these constants, must then
be smaller.

Choose capacities, for example,
\[
                            A_j=2^{j+6}(1+B_j).
\]
The probability that the corresponding upper budget exceeds $A_j$
is at most $2^{-j-6}$.  The small objectives are finitely many
aggregate nonnegative random variables.  Their expected values
can be made smaller than their desired bounds by any prescribed
factor, using the initial affordable tails and the subsequent
absolute errors already described.  Fix these with enough room
that their Markov failure probabilities sum to less than $1/8$.
The sharp affine-error objective is treated in the same way.
The sum of the crude-capacity failure probabilities is also less
than $1/8$.  Thus a positive-probability set of edge jitters
satisfies all these upper-budget inequalities.

\emph{Guards, meshes, and samples.}
Now use these fixed capacities in Proposition~\ref{wl:barrier} to choose
the guard spacings and their buffers.  The desired accuracies include
those used in Subsection~\ref{ha:active-localization}, all local value tests,
and the following sharper requirement on every saturation cube:
\begin{equation}\label{ha:sharp-label-pinning}
                  Y=(S_D(t(q_h)),S_e(s(q_h))).
\end{equation}
Require this pair to be uniformly as close to $Y_0$ as prescribed
relative to the aligned cube side length.  Both labels stay in the
same smooth sector, and the activation
function $H_{a_D}$ is the identity at their resulting radii.
These are physical label requirements after the finite patches
have been selected.  They can be imposed by a sufficiently fine
combined label accuracy and sufficiently close stairs.

Properness bounds the number of guard-label gap exclusions relevant
to a compact feature.  Choose the preliminary guard realization
accuracies relative to these prescribed gap budgets, which may be
smaller than the relative tolerance $\zeta$ to ensure fine stair
motion.  The dyadic sampling belts in
Lemma~\ref{ms:boundary-collar} can meet these prescriptions by
using summably small fractions across scales and requiring every
scale actually used on a compact feature to lie sufficiently far
down the tail.  Only finitely many such scales are used there
once a mesh has been selected.

Choose $d_0$ and the meshes small enough for the guard, high-label
and regional buffers.  The size functions have constant cores and
a common sharp-interior step $H$, as in Lemma~\ref{ms:meshes}.
The push into the cut boundary preserves the already fixed local
bi-Lipschitz constants.  Apply the allowed boundary and volume
jitter.  The uniform upper-budget bounds just proved apply to this
mesh, so choose a jitter satisfying them and all almost-sure ACL
conditions.  Finally choose the central sampling weights after
these edges.  In the finite equal-weight regions take $W/H$ as
small as required.  Split allowed components sufficiently finely
and place intervals $I_j$ of width comparable to $c_*w_j$, with
small sampling margins.  Lemma~\ref{wl:stairs} makes all required
weighted count inequalities hold with simultaneous positive
conditional probability.  Countably many local sampling tests
use summable conditional failure probabilities.  Hence at least
one joint choice of edges and samples satisfies everything.
The displacement of whole-interval and occupied-interval stairs
is charged only to absolute gaps, rounding and slot size; keeping
$c_*>0$ fixed introduces no residual motion error.

In the same conditional selection impose the almost-sure tests
needed to open the carrier: avoid labels of nonsingleton fibre
values, labels attained on nonexact seams along edges, and mesh
vertices; require finite original-carrier hit sets, including in
the exact collar.  These tests use only coarea and null avoidance,
and add no numerical energy constraint.  Their applicability is
verified in the next proof.  Intersecting them with the
positive-probability set just obtained completes the selection.
\end{proof}

\subsection{From sampled counts to the actual wall systems}
\label{ha:wall-transfer}

The sampled estimates refer to the original carrier away from the
boundary adjustment and to the prepared trace family inside its
exact collar.  We now obtain actual walls with these counts.  A
single opening is used for both colours; its equality with the
carrier in the exact collar then allows the boundary preparation
to realize the second set of tests.

\begin{proposition}[Initial walls with the selected budgets]
\label{ha:initial-walls}
The guards, mesh, samples and count bounds selected in
Lemma~\ref{ha:capacities} admit pre-stacks satisfying
Definition~\ref{wl:prestack} and avoiding the selected guards.
Their actual slope majorant $G$ satisfies
\[
                         \int_{U_j}G^p\le A_j
\]
on every buffered barrier chart $U_j$, together with the weak and
sharp estimates of Lemmas~\ref{ha:ambient-tests} and
\ref{ha:sharp-counts}.  The mesh and samples are unchanged.
Consequently the hypotheses of Theorem~\ref{wl:realization} hold.
\end{proposition}

\begin{proof}
\emph{A common opening and the initial walls.}
We now pass from the sampled carrier levels to the initial systems
required by Proposition~\ref{ms:prestacks}; this step precedes routing.
Let $\mathcal W$ be the entire two-colour family of sampled standard
walls in $M_y$.  Each wall is compact and properly embedded, and
there are finitely many samples per feature.  The standard feature
family is locally finite.  Hence $\mathcal W$ is locally finite and
its union is relatively closed in $M_y$.  The nonsingleton values
of the central carrier in Lemma~\ref{qt:central-carrier} are
countable.  Each such value forbids at most finitely many radial
or intermediate meridian levels, and a graph vertex is not an
intermediate meridian level.  These countably many labels are among
the almost-sure sample avoidances in Lemma~\ref{ha:capacities}.
Thus the whole wall union,
not merely its edge-hit values, avoids every nonsingleton value.

For the opening use the \emph{original} carrier $F_*$ and the
actual final mesh edges, before the boundary preparation has
changed the walls.  Away from the exact collar, the already chosen
ACL and coarea tests give finite hit sets on compact localizations.
The nonexact seam set $S$ of Lemma~\ref{qt:central-carrier} has
edge-parameter measure zero by volume jitter.  The relevant label
paths are absolutely continuous, so the images of these parameter
sets have one-dimensional measure zero.  These scalar label sets
and the mesh-vertex labels were also excluded in the same sample
selection.  This gives no hit on $S$ outside the exact regions.
In the exact collar the unmodified carrier labels along the final
edges are locally Lipschitz.  Their additional coarea tests therefore
give finite hits almost surely; these tests require no numerical
energy budget.  The finite feature lists permit all of them in
the selection of Lemma~\ref{ha:capacities}.  A
compact localization meets finitely many edges and sampled walls,
so its total actual hit set is finite.

Lemma~\ref{qt:edge-matching} therefore supplies \emph{one common}
opening homeomorphism $f_{\mathrm o}:\Omega\to\Lambda$ for the
entire two-colour wall union, with exactly its original carrier
hits on every mesh edge.  Choose it sufficiently finely close to
$F_*$ for the prescribed guard-gap buffers.  It equals $F_*=h_*$
on an exact neighborhood containing all supports of the boundary
preparation, as well as on the other prescribed exact data, and
maps $M_x$ onto $M_y$ with the fixed boundary correspondence.
Pull back both standard families by this same $f_{\mathrm o}$.
They are compact properly embedded locally flat systems of the
required individual types, with noncontractible annular cores.

Now apply the common joint boundary adjustment of
Lemma~\ref{ms:boundary-collar}: the warp and chord preparation,
its mirrored continuation inside the cut, and the depth stretch.
This operation is supported entirely in the exact neighborhood
where $f_{\mathrm o}=F_*=h_*$.  Therefore the resulting collar
walls and their actual edge counts are precisely the prepared
data whose weighted coarea bounds were used above.  Outside that
neighborhood their edge counts are the original carrier counts,
preserved by the opening.  The interval-ribbon adaptation keeps
the central sample chords fixed.  Consequently all previously
selected weighted and directional upper counts apply to these
actual initial systems.  No sampling estimate on a
sample-dependent pullback of an edge is being used.

Apply Lemma~\ref{wl:relative-pl} separately to each colour,
retaining the smaller prescribed PL boundary product collars.
Its hypotheses hold by the finite isolated edge hits, vertex
avoidance and the exact prepared boundary data.  Choose its fine
position tolerances within the guard buffers.  It gives the
required PL systems and flat transverse interior crossing germs,
without increasing any individual wall--edge count.  Apply
Proposition~\ref{wl:normalization} separately to these two systems,
retaining their boundary traces, topologies and coorientations,
and again without increasing an individual edge count.  Neither
application requires a relative isotopy fixing the other colour.
Proposition~\ref{ms:prestacks} now supplies the pre-stacks.
No individual edge count has increased.  The same count-based
slope upper budgets selected in Lemma~\ref{ha:capacities} therefore
bound the actual integrals $\int_{U_j}G^p$ by $A_j$ and give the
stated weak and sharp estimates.  Their guard avoidance follows from the
reserved label gaps and the chosen finer mesh: each final normal
disk lies in a tetrahedron incident to an original possible hit,
and placements, thickening and snapping stay in its prescribed
mesh buffers.  The preliminary opening and PL changes were chosen
within the same positive buffers.  Thus this execution establishes
the pre-stack hypotheses of Theorem~\ref{wl:realization} without
changing the previously fixed capacities or the guard order.

\emph{Slopes on lower-dimensional interfaces.}
For the rectifiable-path use of the barrier, define $G$ on mesh
interfaces and exceptional collar-boundary strata to include the
larger adjacent actual local collar slopes.  These strata have
ambient volume zero and locally finite prescriptions, so this
changes none of the volume budgets.  The pre-parameters are
Lipschitz on their compact product collars; restricted path coarea
therefore uses these pointwise upper slopes.  On open pieces the
volume bounds are exactly those already proved.  In particular,
no estimate for an inverse distance across a gap at a manipulated
boundary is needed in the integral capacity.
\end{proof}

\begin{remark}\label{ha:scale-order}
The order established above can be recorded without a forward
parameter dependency.  Fix $p$ and the desired power error; choose
$\gamma,\eta$, then $\zeta,c_*$ and the ordinary weak constant.
Choose affordable rank tails and finite jet truncations, then the
full-rank shell tests.  Construct the carrier and fix its regional
and weighted marked-collar budgets.  Fix the polar data, aspect
localizations, full-cell layer widths, activation radii and exact
central carrier.  Fix compact sharp-template and basis bounds,
then level tolerances, chart/jitter fractions and smaller
power-mean tests; select the finite sharp patches.  Choose the paid
transition regions and the crude capacities.  Only then choose
guard accuracies and spacings.  Finally take the boundary widths
and meshes fine enough for all buffers, select a permissible
jitter, and sample with sufficiently small weights.  The expectation
bounds are uniform in these last resolutions, which is why the
capacities may precede the guards.  Strict parameters, marked-cube
collapses and the final smoothing are chosen afterward.
\end{remark}

\subsection{Saturation and physical gradient accuracy}
\label{ha:saturation-subsection}

Apply the wall construction to the pre-stacks of
Proposition~\ref{ha:initial-walls}, with the capacities and samples
selected in Lemma~\ref{ha:capacities}.  Theorem~\ref{wl:realization} gives the
locally bi-Lipschitz $h$, the actual occupied intervals and the
label pinning.  We estimate $P_0Th$, not the derivative of the
qualitative product and ball extension used to construct $h$.

\begin{lemma}[Saturation and physical gradient accuracy]\label{ha:saturation}
On the union of the saturation interiors,
\begin{equation}\label{ha:physical-saturation}
 \int |D(P_0Th)-Df|^p
       \le C_p(\eta+\gamma)(1+E)
                         +\text{freely small errors}.
\end{equation}
The constant in the displayed leading term is independent of the
compact jet and chart bounds.
\end{lemma}

\begin{proof}
In a saturation cube $Q$ in aligned coordinates put
\[
                  \widetilde Y=Y\circ g',\qquad D_0=BV.
\]
Proposition~\ref{ms:sharp-contract} and
Proposition~\ref{wl:parameter-cost} give, off the dirty portions,
\begin{equation}\label{ha:Frobenius}
 |D_y\widetilde Y|\le
 \begin{cases}
   \sqrt2+C\eta,&\rank D_0=2,\\
   1+C\eta,&\rank D_0=1.
 \end{cases}
\end{equation}
At an essential both-active crossing, the two scalar gradients
are separately near-unit orthogonal axis gradients.  On a switched
diagonal only one output scalar is active and its derivative is
their sum or difference, with the same squared-norm bound.  In
parallel mode the opponent strips in this deep flat region are
disjoint.  Strip width $a_1$, flat slopes and occupied-root speeds
contribute $O(\eta)$ after the subsequent smaller choices.  Thus
arbitrary partial switches do not increase the upper bound in
\eqref{ha:Frobenius}.  By Lemma~\ref{ha:sharp-counts}, the evaluated
$p$-integrals on dirty portions are arbitrarily small in aggregate,
with any fixed physical conversion weights.

The mean gradient has the required affine value:
\begin{equation}\label{ha:mean-gradient}
                \left|\frac{1}{|Q|}\int_Q D_y\widetilde Y-D_0\right|
                      \quad\hbox{is arbitrarily small}.
\end{equation}
For $Y_0\circ g'$ this follows from the enlarged power-mean tests,
the bounded Jacobian factors and the chart/jitter error.  For the
difference, integration on the boundary of the cube gives
\[
 \left|\frac{1}{|Q|}\int_Q D_y\bigl((Y-Y_0)\circ g'\bigr)\right|
       \le \frac{C}{\operatorname{side}Q}
                       \|(Y-Y_0)\circ g'\|_{L^\infty(Q)}.
\]
This is valid for the continuous Sobolev functions at issue by
the trace formula; the right side is small by
\eqref{ha:sharp-label-pinning}.  The cube remains compactly in one
smooth sector.

We use the elementary uniform-convexity inequality
\begin{equation}\label{ha:convexity}
 c_p|Z-D_0|^p\le |Z|^p-|D_0|^p
             -p|D_0|^{p-2}D_0:(Z-D_0),\qquad p>2.
\end{equation}
One proof integrates the Hessian of $|X|^p$ along the segment
$D_0+t(Z-D_0)$.  Its quadratic form is at least
$p|X|^{p-2}|Z-D_0|^2$.  After scaling $|Z-D_0|$ to one, an interval
of fixed positive length in $0\le t\le3/4$ is a fixed positive
distance from the possible zero of the projection of that segment
onto $Z-D_0$.  The integrated Hessian is therefore bounded below
by a positive constant depending only on $p$.  This proves
\eqref{ha:convexity} uniformly in both matrices.

In rank two, $|D_0|=\sqrt2$.  In rank one,
\eqref{ha:rank-one-basis} and \eqref{ha:b-small} imply
$1-C\gamma\le|D_0|\le1$.  Integrate
\eqref{ha:convexity} with $Z=D_y\widetilde Y$.
Equation~\eqref{ha:Frobenius} bounds its energy on the clean set;
the dirty energy is already small.  The linear term is controlled
by \eqref{ha:mean-gradient}, since $|D_0|\le\sqrt2$.  We obtain
\begin{equation}\label{ha:normalized-saturation}
 \int_Q|D_y\widetilde Y-D_0|^p
      \le C_p(\eta+\gamma)|Q|
                   +\text{mean and dirty errors}.
\end{equation}

It remains to verify that the physical conversion does not multiply
the leading error by a large jet condition number.  On this cube
$P_0q_h=\mathcal R(Y)$, with the smooth sector reconstruction
\[
       \mathcal R(t,s)=d_D+\exp(t/L_D)(\xi_e(s)-d_D).
\]
There is no central activation transition.  Since $s$ is arclength
and $L_D$ controls the polygon diameter, its two derivative columns
are bounded by $Cr$; this bound is independent of the thin
rectangle aspect.  Require the chart error to satisfy
$\|(Dg'-V)V^{-1}\|_{\mathrm{op}}\ll1$.  Then
\begin{equation}\label{ha:physical-multiplier}
 \|D\mathcal R(\widetilde Y)\|_{\mathrm{op}}
       \|(Dg')^{-1}\|_{\mathrm{op}}
       \le Cr(z)\|V^{-1}\|_{\mathrm{op}}
       \le C|Df(z)|,
\end{equation}
and $|\det Dg'|=(1+o(1))|\det V|$ relatively.  All these
requirements are possible after the finite basis bounds.

Split the physical derivative into the term containing
$D_y\widetilde Y-D_0$ and the comparison term
\[
             D\mathcal R(\widetilde Y)D_0(Dg')^{-1}.
\]
The latter differs as little as prescribed from $Dq_0(z)$, by the
label, patch and chart tests.  At the central data the identity
$D\mathcal R(Y_0(z))B=Dq_0(z)$ holds.  Now use
\eqref{ha:relative-calibration} and the affine-gradient tests on
$Df$ to compare further with $Df(x)$.  The mean, dirty and comparison
errors may use every finite conversion factor because they were
chosen after those factors.  For the leading term, however,
\eqref{ha:physical-multiplier} and
\eqref{ha:normalized-saturation} give exactly
\[
          C_p(\eta+\gamma)
                      |Df(z)|^p|\det V|\,|Q|.
\]
The selected outer patches are disjoint, their affine-gradient
tests are accurate, and their chart Jacobians are relatively close
to $|\det V|$.  Hence
\[
                  \sum_Q|Df(z_Q)|^p|\det V_Q|\,|Q|
                         \le C(E+1).
\]
Summing proves \eqref{ha:physical-saturation}.  The flattening angle
and central-fraction losses in \eqref{ha:relative-calibration} were
freely chosen small relative to this same finite energy.
\end{proof}

\subsection{Strict maps, marked interiors and completion}
\label{ha:completion}

\begin{proof}[Proof of Theorem~\ref{ha:approximation}]
Carry out the preceding choices with sufficient room in each target
error.  Let $u=P_0Th$ be the nonstrict map.  On saturation interiors
it satisfies Lemma~\ref{ha:saturation}.  On the exterior residual,
full-rank shells and weighted marked collars, the derivative costs
are arbitrarily small by Lemma~\ref{ha:ambient-tests} and the
simultaneous choice and transfer in Lemma~\ref{ha:capacities} and
Proposition~\ref{ha:initial-walls}.  The $|Df|^p$
integral on the residual is also small: the positive deficient ranks
are captured by the saturation interiors up to the affordable
losses, the full-rank part is captured by its inner cubes up to the
chosen tail, and the rank-zero derivative vanishes.  The corresponding
$|Df|^p$ costs on the correction shells and enlarged marked cubes
were fixed small at the initial stage.

Use Proposition~\ref{wl:strict} and Lemma~\ref{qt:strict} to
replace $u$ by
\[
                            k=P^\epsilon T_\delta h.
\]
The symbols $\epsilon$ and $\delta$ here denote the strict target
parameters, not the slot loss.  Choose them locally with summably
small derivative and value errors on every prescribed test, including
the finite full-rank tests and the marked-collar weights.  The stated
tangential compatibility on faces and edges ensures this derivative
convergence; at vertex levels the nonstrict derivative is zero.
The map $k$ is a locally bi-Lipschitz homeomorphism of $\Omega$ onto
$\Lambda$.  We impose no global energy bound on its still exempt
marked interiors at this stage.

If $2<p<3$, apply Lemma~\ref{hp:hide-cubes}, using
\eqref{ha:collar-budget} and its transferred weighted volume tests
for $k$.  Choose the rescaled boundary in a slightly outer collar
so that the entire exempt cube is enclosed.  We recall why this
step can make its energy small despite an arbitrarily large fixed
interior Lipschitz constant.  In cube polar coordinates, the source
self-homeomorphism sends $[0,\alpha]$ linearly to $[0,1-\rho]$ and
$[\alpha,1]$ to $[1-\rho,1]$, fixing the boundary.  The inner energy
is bounded by a fixed compact constant times $\alpha^{3-p}$ and
can be made arbitrarily small after $k$ has been chosen.  For fixed
$\alpha$, take $\rho\downarrow0$ at the chosen radial Lebesgue
boundary.  The outer energy is bounded by a constant times
\[
 \ell_U\int_{\text{chosen facets}}|Dk|^p
                     \int_\alpha^1 t^{2-p}\,dt,
 \qquad \int_\alpha^1 t^{2-p}\,dt\le \frac{1}{3-p}.
\]
Slice selection bounds the first factor by the weighted collar
volume test.  The factor depending on $p<3$ is fixed before those
budgets are prescribed.  Use summably small choices for the marked
cubes, also including any finite full-rank shell factors.  The changes
avoid all saturation patches.  Their value error is small because
$f$ has the prescribed small oscillation on each enlarged cube and
$k$ already meets the fine value test there.  Denote the resulting
locally bi-Lipschitz homeomorphism by $k_1$.  If $p\ge3$, put
$k_1=k$.

The map $k_1$ now belongs to global $W^{1,p}$.  The residual,
saturation, shell and marked-interior derivative costs have summable
bounds.  The only possibly unpaid full-rank interiors lie in a
finite compact subset of $\Omega$, where local bi-Lipschitzness
gives finite $p$-energy.  Finally its values lie in the bounded
$\Lambda$, and $\Omega$ has finite measure.  This proves actual
membership, rather than merely local membership, before invoking
the smoothing Theorem.

Apply Theorem~\ref{sm:smoothing} to $k_1$, keeping room in the
energy and value tests, in particular the normalized full-rank
sup tests.  Perform the finitely many smooth cubical jet corrections
of Proposition~\ref{hp:full-rank-correction}.  The map equals $b_i$
on the exact inner cubes, where the error is
\eqref{ha:affine-full-tests}.  Its changed shells have arbitrarily
small cost by the buffered budgets including the fixed jet factors.
Their supports avoid the saturation patches.  On the complement,
the sharp estimate remains in force and the residual small-energy
estimates for both maps control the gradient difference.  Smooth
once more by Theorem~\ref{sm:smoothing}, with an arbitrarily small
additional $W^{1,p}$ error.

All operations have the required fixed target.  The strict target
maps are self-homeomorphisms of $\Lambda$; the wall realization
maps onto $\Lambda$; the marked-cube operations are source
self-homeomorphisms fixed at their boundaries; the cubical
corrections preserve their previous images; and the smoothing
Theorem preserves the source and target.  Local finiteness and the
fine properness tests in the preceding constructions apply on the
open domains throughout.
Thus surjectivity holds for every individual approximant, rather
than only for a limiting image.  Its everywhere invertible smooth
differential supplies local smooth inverses, which bijectivity
assembles into a smooth inverse on $\Lambda$.

We may also make the value power error arbitrarily small.  One
explicit way to pass from the fine local tests to this assertion is
to choose a compact subset of $\Omega$ whose complement has
arbitrarily small measure.  On the compact impose arbitrarily small
uniform value error through the preliminary constructions,
corrections and smoothing.  On the complement use the fixed bound
on the diameter of $\Lambda$.  Full-rank correction supports and
marked cubes were chosen with the additional small oscillation
requirements needed for their localized motions.  Thus none of the
last operations obstructs this value estimate.

Combining the estimates gives a final derivative error bounded by
\[
                    C_p(\eta+\gamma)(1+E)+o(1),
\]
where $o(1)$ denotes the later freely prescribed errors, after all
of their multipliers are fixed.  First choose $\eta,\gamma$ so that
the displayed leading term is smaller than the desired error share;
then choose the tails against $C_{\mathrm w}$ and the successive
absolute errors in Remark~\ref{ha:scale-order}.  The value error
and the two smoothing errors receive the remaining shares.  This
gives the prescribed bound $\eps$.  Apply the construction with
$\eps=2^{-j}$ and undo the initial reflection if needed.
The resulting sequence satisfies Equation~\eqref{main:convergence},
with each approximant in global $W^{1,p}$ and onto the same fixed target.
\end{proof}
